\documentclass[11pt]{article}
\usepackage{fontspec}
\usepackage[french]{babel}
\usepackage{amsmath}
\usepackage{amssymb}
\usepackage{amsthm}
\usepackage{mathrsfs}
\usepackage{tikz-cd}
\usepackage{enumitem}
\usetikzlibrary{babel}

\newcommand{\printedpage}[1]{\marginpar{\footnotesize p.~#1}}
\newcommand{\msp}[1]{\marginpar{\footnotesize [#1]}}
\newcommand{\Ind}{\operatorname{Ind}}
\newcommand{\Ob}{\operatorname{Ob}}
\newcommand{\Ord}{\operatorname{Ord}}
\newcommand{\card}{\operatorname{card}}
\newcommand{\cardred}{\mathop{\textup{card\,réd}}\nolimits}
\newcommand{\cercle}[1]{\tikz[baseline=(c.base)]\node[draw,circle,inner sep=1pt](c){\textup{#1}};}

\begin{document}
\setcounter{footnote}{143}

\printedpage{129}
% head: A. Grothendieck : Dérivateurs, Ch. XVIII, Édition des Universités de Montpellier II et Paris VII
% foot: Version 17 décembre 2025

\msp{page 340}
à l'exception de a$'$ $\Longrightarrow$ a, qui est un peu plus simple dans le cas actuel. On va la dégager sous forme de lemme un peu plus général. (\textbf{N.B.} La démonstration de a$'$ $\Longrightarrow$ a dans 4.13.15 avait été omise comme “soritale”\,!, cf. page 337\,\dots)

\textbf{Lemme 4.13.19.} \textit{Soient $\mathcal{M}$ une catégorie accessible dans $\mathcal{M}$, $\pi$ un cardinal utile pour $\mathcal{M}$ et tel que $\underline{S}$ soit $\pi$-accessible, i.e. $\underline{S}$ est l'image essentielle de $\Ind_\pi(\underline{S}(\pi))$ par le foncteur pleinement fidèle canonique $\Ind_\pi(\underline{S}(\pi))\longrightarrow\mathcal{M}$, où $\underline{S}(\pi)=\underline{S}\cap\mathcal{M}_\pi$. On suppose qu'il existe un ordinal limite $\alpha$, tel que $\card\alpha\leqslant\pi$ et tel que $\mathcal{M}_\pi$ et $\underline{S}(\pi)$ $(=\underline{S}_\pi)$ soient stables par $\varinjlim$ de type $I_\alpha$ soient stables par $\varinjlim$ de type $I_\alpha$, et que le foncteur d'inclusion $\underline{S}(\pi)\longrightarrow\mathcal{M}_\pi$ y commute. Soit $\Phi$ un ensemble de catégories telles que $\forall\ I\in\Phi$, on ait $\cardred I\leqslant\pi$, et supposons que $\mathcal{M}_\pi$ soit stable par $\varinjlim$ de type $\Phi$, et que $\underline{S}(\pi)$ soit stable dans $\mathcal{M}_\pi$ par lesdites limites (ou ce qui revient au même, qu'il soit stable en lui-même par ces $\varinjlim$\,, et que le foncteur d'inclusion y commute). Alors $\mathcal{M}$ et $\underline{S}$ sont stables par $\varinjlim$ de type $\Phi$, et le foncteur d'inclusion $\underline{S}\longrightarrow\mathcal{M}$ y commute, i.e. $\mathcal{M}$ est stable par ces limites, et $\underline{S}$ stable dans $\mathcal{M}$ par ces mêmes limites.} % [sic] « soient stables par lim de type I_alpha » printed twice -> once, repetition

\msp{page 341}
\textsc{Démonstration}. On applique 4.11.1 (p. 267-268) séparément à $\mathcal{M}\simeq\Ind_\pi\mathcal{M}_\pi$ et à $\underline{S}\simeq\Ind_\pi\underline{S}(\pi)$, pour conclure que $\mathcal{M}$ et $\underline{S}$ sont stables séparément par $\varinjlim$ de type $\Phi$. Il reste à voir que le foncteur $\underline{S}\longrightarrow\mathcal{M}$ y commute. Cela se voit comme dans la démonstration de 4.11.1, on obtient plus généralement ceci :

\textbf{Lemme 4.13.20} (cor. à 4.11.1). \textit{Soient $f_0 : C\longrightarrow C'$ un foncteur entre catégories équivalentes à des petites catégories, $\pi$ un cardinal infini, $f:\mathcal{M}\stackrel{\textup{déf}}{=}\Ind_\pi C\longrightarrow\mathcal{M}'\stackrel{\textup{déf}}{=}\Ind_\pi C'$ l'extension canonique $\Ind_\pi(f_0)$ de $f_0$. Soient $\Phi$ et $\alpha$ comme dans 4.11.1, et supposons que $I_\alpha\in\Phi$ et que $f_0$ commute aux $\varinjlim$ de type $\Phi$. Alors $f$ commute à ces mêmes limites inductives.} (\footnote{Il vaut mieux prendre des foncteurs $f_0 : C\longrightarrow\mathcal{N}$, et $\mathcal{N}$ [?] que $f_0$ commute aux $\varinjlim$ de type $\Phi$, donc il en est de même de son extension $\Ind_\pi(C,\mathcal{N})\longrightarrow\mathcal{N}$. (\textbf{N.B.} On suppose que $\mathcal{N}$ satisfait $\mathrm{L}_\pi$.)}).

\textbf{Corollaire 4.13.21.} \textit{Sous les conditions de 4.13.20, supposons que $\Phi$ contienne les catégories de cardinal $\leqslant\pi$ (resp. les catégories filtrantes de cardinal $\leqslant\pi$, resp. les catégories filtrantes grandes devant $\pi_0$ de cardinal $\leqslant\pi$, resp. les catégories discrètes de cardinal $\leqslant\pi$). Alors $f$ commute aux petites $\varinjlim$ (resp. aux petites $\varinjlim$ filtrantes, resp. aux petites $\varinjlim$ filtrantes grandes devant $\pi_0$, resp. aux petites sommes). (Réciproque claire et “laissée au lecteur”.)}

Se voit comme 4.4.14 et 4.4.14.1 (p. 97, 98).

\msp{page 342}
Voici enfin, il me semble, la forme définitive (\,?) de 4.13.15, 4.13.18.

\textbf{Théorème 4.13.22} (\footnote{cf. variante 4.13.36, page 358.}). \textit{Soient $\mathcal{M}$ une catégorie accessible, $\pi$ un cardinal infini, utile pour $\mathcal{M}$, et $\Phi$ un ensemble de catégories satisfaisant les deux conditions :}
\begin{itemize}[leftmargin=2.5em]
\item[\textit{1°)}] $J\in\Phi$ \textit{implique} $\cardred J\leqslant\pi$.
\item[\textit{2°)}] \textit{Il existe un} ordinal limite $\alpha_0$ \textit{avec} $I_{\alpha_0}\in\Phi$.
\end{itemize}
\textit{On suppose $\mathcal{M}$ stable par $\varinjlim$ de type $\Phi$. Soient $\underline{S}$ une sous-catégorie} strictement pleine \textit{de $\mathcal{M}$, et $\underline{S}(\pi)=\underline{S}\cap\mathcal{M}_\pi$. Désignons par $L(\varphi)$ la condition sur $\underline{S}(\pi)$ : $\underline{S}(\pi)$ est stable dans} % [sic] L(\varphi) -> L(\Phi), the condition is later written L(\Phi)
\printedpage{130}
% head: A. Grothendieck : Dérivateurs, Ch. XVIII, Édition des Universités de Montpellier II et Paris VII
% foot: Version 17 décembre 2025
\textit{$\mathcal{M}$ (où ce qui revient au même, dans $\mathcal{M}_\pi$) par $\varinjlim$ de type $\Phi$. Ceci posé, les conditions suivantes sur $\underline{S}$ sont équivalentes :} % [sic] où ce qui revient -> ou ce qui revient, conjunction "ou"
\begin{itemize}[leftmargin=2.5em]
\item[\cercle{a}] \textit{$\underline{S}$ est $\pi$-accessible dans $\mathcal{M}$, et stable dans $\mathcal{M}$ par $\varinjlim$ de type $\Phi$.}
\item[a$'$)] \textit{$\underline{S}$ est $\pi$-accessible dans $\mathcal{M}$, et $\underline{S}(\pi)$ satisfait à $L(\Phi)$.}
\end{itemize}
\msp{page 343}
\begin{itemize}[leftmargin=2.5em]
\item[b)] \textit{$\underline{S}(\pi)$ satisfait à $L(\Phi)$, et pour tout $x\in\Ob\mathcal{M}$, on a ceci : $x\in\Ob\underline{S}$ si et seulement si la sous-catégorie strictement pleine $\underline{S}(\pi)/x$ de $\mathcal{M}(\pi)/x$ est cofinale dans $\mathcal{M}(\pi)/x$} \texttt{[lire} $\mathcal{M}_\pi/x$\texttt{]}.
\item[\cercle{c}] \textit{$\underline{S}(\pi)$ satisfait à $L(\Phi)$, et pour tout $x\in\Ob\mathcal{M}$ et toute représentation $x\simeq\varinjlim\limits_{i\in I}x_i$ qui est $\Phi$-$\pi$-adaptée (cf. définition plus bas), on a ceci : $x\in\Ob\underline{S}$ si et seulement si $I_{\underline{S}}$ est cofinale dans $I$, où $I_{\underline{S}}$ est la sous-catégorie strictement pleine de $I$ formée des $i\in\Ob I$ tels que $x_i\in\Ob\underline{S}$ (ou encore, $x_i\in\Ob\underline{S}(\pi)$).} \textbf{N.B.} \textit{On dit qu'un système inductif $(x_i)_{i\in I}$ dans $\mathcal{M}$ est $\Phi$-$\pi$-adapté si}
\begin{itemize}[leftmargin=2em]
\item[\textit{1)}] \textit{$I$ équivalente à une petite catégorie, et $I$ grande devant $\pi$,}
\item[\textit{2)}] \textit{les $\varinjlim$ de type $\Phi$ existent dans $I$,}
\item[\textit{3)}] \textit{le foncteur $i\longmapsto x_i$ de $I$ dans $\mathcal{M}$ commute aux $\varinjlim$ de type $\Phi$, et enfin}
\item[\textit{4)}] \textit{les $x_i\in\Ob\mathcal{M}_\pi$.}
\end{itemize}
\item[c$'$)] \textit{(Si $\pi^{\pi_0}=\pi$, où $\pi_0$ est un cardinal $\leqslant\pi$ utile pour $\mathcal{M}$, et si $J\in\Phi$ implique que $J$ est grande devant $\pi_0$.) Comme c), avec $I$} ordonnée.
\end{itemize}
(\footnote{\textbf{N.B.} Quand les conditions restrictives ci-contre \texttt{[parenthèse dans c$'$)]} ne sont pas satisfaites, alors a, a$'$, b, c sont encore équivalentes, et impliquent c$'$, mais ne sont peut-être pas équivalentes à c$'$).},\footnote{\textbf{N.B.} Il faudrait réénoncer les conditions annexes [?] (Si \dots), cf. p. 359-360.}).

\textbf{Définition 4.13.23} (\footnote{\textbf{N.B.} Il y aurait intérêt à remonter cette définition (\textit{sans} b), dès avant 4.13.15.}). Soient $\mathcal{M}$ une catégorie accessible, $\pi$ un cardinal, $\Phi$ un ensemble de catégories, satisfaisant les deux conditions :
\begin{itemize}[leftmargin=2.5em]
\item[a)] $J\in\Phi\Longrightarrow\cardred J\leqslant\pi$.
\item[b)] Il existe un ordinal limite $\alpha_0$, tel que $I_{\alpha_0}\in\Phi$ (\footnote{Si c'est pour poser une définition, on peut se passer de la condition b).}).
\end{itemize}
\msp{page 344}
Soit $S$ une partie de $\Ob\mathcal{M}$, $\underline{S}$ la sous-catégorie pleine qu'elle définit. On dit que $S$ est une \textit{partie $\Phi$-$\pi$-accessible de} $\Ob\mathcal{M}$, ou que $\underline{S}$ est une \textit{sous-catégorie $\Phi$-$\pi$-accessible dans} $\mathcal{M}$, si les conditions suivantes sont satisfaites.
\begin{itemize}[leftmargin=2.5em]
\item[1°)] $\pi$ est utile pour $\mathcal{M}$, $\mathcal{M}$ est stable par $\varinjlim$ de type $\Phi$ (\footnote{Condition préliminaire sur $\mathcal{M}$.}).
\item[2°)] $\underline{S}$ est une sous-catégorie $\pi$-accessible dans $\mathcal{M}$.
\item[3°)] $\underline{S}$ est stable dans $\mathcal{M}$ par $\varinjlim$ de type $\Phi$.
\end{itemize}
(\footnote{\textbf{N.B.} Quand $\Phi=\{I_{\alpha_0}\}$, on dit aussi que $\underline{S}$ est $\alpha_0$-$\pi$-accessible dans $\mathcal{M}$. Si $\Phi$ est formé des catégories de cardinal $\leqslant\pi$ et grandes devant $\pi_0$, on dit que $\underline{S}$ est $\pi_0$-$\pi$-accessible dans $\mathcal{M}$.}).

\textbf{N.B.} Quand la condition préliminaire 1°) \textit{sur} $\mathcal{M}$ est satisfaite, la validité des conditions 2°) 3°) équivaut donc à la condition que $\underline{S}$ soit \textit{strictement} pleine dans $\mathcal{M}$, et que de
\printedpage{131}
% head: A. Grothendieck : Dérivateurs, Ch. XVIII, Édition des Universités de Montpellier II et Paris VII
% foot: Version 17 décembre 2025
plus $\underline{S}$ satisfasse la condition a) du théorème 4.13.22. D'autre part, cette condition est équivalente à chacune des conditions a$'$, b, c de ce même théorème, et aussi à la dernière condition c$'$, sous réserve qu'on ait $\pi^{\pi_0}=\pi$ (pour un cardinal $\pi_0<\pi$ utile pour $\mathcal{M}$) et que $J\in\Phi$ implique que $J$ est grande devant $\pi_0$.

Conceptuellement, c'est la définition donnée dans 4.13.23, i.e. la condition \cercle{a} de 4.13.22, qui paraît la plus claire, ou sa variante immédiate a$'$. Mais du point de vue pratique, du travail avec les parties accessibles, c'est par propriété \cercle{c} qui paraît la plus commode. Le poids principal du théorème 4.13.22 est bien
\msp{page 345}
l'implication a $\Longrightarrow$ c (via a $\Longrightarrow$ a$'$ $\overset{!}{\Longrightarrow}$ c). Donc il faut surtout voir ce théorème comme donnant un critère remarquable d'appartenance à $\Ob\underline{S}$ d'un objet $x$ de $\mathcal{M}$ qui est présenté par un système inductif $\Phi$-$\pi$-adapté :
\[ x\ \in\ \Ob\underline{S}\qquad\quad\Longleftrightarrow\qquad\quad I_{\underline{S}}\ \text{cofinal}\ , \tag*{$(*)$} \]
et dans cette équivalence, c'est l'implication $\Longrightarrow$ qui est la partie non soritale, et remarquable.

\textbf{Proposition 4.13.24.} \textit{Soient $\mathcal{M}$, $\pi$, $\Phi$ comme dans 4.13.23} (\footnote{\textbf{N.B.} La condition $I_{\alpha_0}\in\Phi$ n'est pas utile.})\textit{, on suppose que la condition 1° de 4.13.23 est satisfaite ($\pi$ utile pour $\mathcal{M}$, $\mathcal{M}$ stable par $\varinjlim$ de type $\Phi$). Soit ${\Sigma'}^{\Phi}_{\pi}$ l'ensemble des sous-catégories $\underline{S}$ de $\mathcal{M}$ qui sont $\Phi$-$\pi$-accessibles dans $\mathcal{M}$, et $\Sigma^{\Phi}_{\pi}$ l'ensemble des sous-catégories $\underline{S}(\pi)$ de $\mathcal{M}_\pi$ qui sont strictement pleines, et stables dans $\mathcal{M}_\pi$ par $\varinjlim$ de type $\Phi$. (Ces sous-catégories sont forcément karoubiennes, car stables par $\varinjlim$ de type $I_{\alpha_0}$, donc aussi par facteurs directs.) Alors les bijections (4.13.4.1) (page 288) induisent des bijections inverses l'une à l'autre}
\[
\begin{tikzcd}
\Sigma^{\Phi}_{\pi} \arrow[r, shift left] & {\Sigma'}^{\Phi}_{\pi} \arrow[l, shift left]
\end{tikzcd}\ .
\]

\msp{page 346}
% marge : trait manuscrit (reproduit en image) en regard de la Proposition 4.13.25
\textbf{Proposition 4.13.25.} \textit{Soient $\mathcal{M}$, $\pi$, $\Phi$ comme dans 4.13.23, avec $I_{\alpha_\pi}\in\Phi$ ($\alpha_\pi$ étant le plus petit ordinal de cardinal $\Phi$} \texttt{[plutôt} $\pi$\texttt{]}\textit{). Alors l'ensemble ${\Sigma'}^{\Phi}_{\pi}$ des sous-catégories de $\mathcal{M}$ $\Phi$-$\pi$-accessibles dans $\mathcal{M}$ est stable par intersection de familles de cardinal $\leqslant\pi$} (\footnote{Si $\pi$ n'est pas cardinal limite. Sinon, il faut dire : familles de cardinal $<\pi$. Cf. 4.13.10.3, où je donne un énoncé rectifié légèrement plus général.}).

(\textbf{N.B.} La bijection $\Sigma^{\Phi}_{\pi}\xrightarrow{\ \sim\ }{\Sigma'}^{\Phi}_{\pi}$ de 4.13.24 respecte les relations d'ordre naturelles d'inclusion, donc transforme intersections en intersections.)

Se voit comme 4.13.7 (qui est contenu dans l'énoncé actuel), en utilisant le lemme 4.13.10.

\textbf{4.13.26.} Soient $\mathcal{M}$, $\pi$, $\Phi$ comme dans 4.13.23, et soit $\underline{S}$ une sous-catégorie \textit{strictement pleine} de $\mathcal{M}$, $f:\underline{S}\longrightarrow\mathcal{M}$ le foncteur d'inclusion. La condition que $\underline{S}$ soit $\Phi$-$\pi$-accessible dans $\mathcal{M}$ se décompose en :
\begin{itemize}[leftmargin=2.5em]
\item[1)] $\pi$ est utile pour $\mathcal{M}$ \textit{et} pour $\underline{S}$.
\item[2)] $\mathcal{M}$ et $\underline{S}$ sont stables par $\varinjlim$ de type $\Phi$.
\item[3)] $f$ commute auxdites limites.
\item[4)] $f$ est $\pi$-accessible.
\item[5)] $f(\underline{S}_\pi)\subset\mathcal{M}_\pi$.
\end{itemize}

\printedpage{132}
% head: A. Grothendieck : Dérivateurs, Ch. XVIII, Édition des Universités de Montpellier II et Paris VII
% foot: Version 17 décembre 2025
\msp{page 347}
Soit alors $\pi'>\pi$ un autre cardinal, quand peut-on assurer que si $\underline{S}$ est $\Phi$-$\pi$-accessible dans $\mathcal{M}$, il soit aussi $\Phi$-$\pi'$-accessible\,? Les conditions 2) 3) 4) sont visiblement stables par passage de $\pi$ à $\pi'$. Restent les conditions 1) et 5). On sait que si $\pi'^{\pi}=\pi'$, alors la condition 1) reste valable pour $\pi'$. Plus généralement, si $\pi_0$ est un cardinal $\leqslant\pi$ tel que $\pi_0$ soit utile pour $\mathcal{M}$ et pour $\underline{S}$, alors si $\pi'>\pi_0$ est tel que $\pi'^{\pi_0}=\pi'$, alors $\pi'$ est aussi utile pour $\mathcal{M}$ et pour $\underline{S}$. Si on admet l'HCG (hypothèse du continu généralisée), il en est ainsi pour tout $\pi'>\pi_0$, donc pour tout $\pi'>\pi$.

Reste la condition 5), pour laquelle une réponse est donnée par 4.12.1 (p. 273). Si $\pi_0\leqslant\pi$ est choisi de façon que $f$ soit $\pi_0$-accessible (p.ex. $\pi_0=\pi$), et si
\[
\begin{aligned}
\pi'_1 &= \cardred\underline{S}_{\pi_0} \qquad\quad (\,\leqslant\cardred\underline{S}(\pi)\,)\\
\pi''_1 &= \sup(\pi'_1,\pi)\ ,
\end{aligned}
\]
on voit que si $\pi'\geqslant\pi''_1$ (i.e. $\pi'\geqslant\pi'_1,\pi$) est tel que $\pi'^{\pi_0}=\pi'$, alors on a bien
\[ f(\underline{S}_{\pi'})\ \subset\ \mathcal{M}_{\pi'}\ . \]

\msp{page 348}
Si donc de plus $\pi'>\pi\geqslant\pi_0$ (comme supposé au début), et si on suppose l'HCG, alors la condition $\pi'^{\pi_0}=\pi'$ est automatique, et il reste la seule condition $\pi'\geqslant\pi''_1$, i.e. en définitive $\pi'\geqslant\cardred\underline{S}_{\pi_0}$. Pour résumer :

\textbf{Proposition 4.13.27.} \textit{Sous les conditions de 4.13.23 (sans condition b) sur $\Phi$), soit $\underline{S}$ une sous-catégorie de $\mathcal{M}$ qui est $\Phi$-$\pi$-adaptée. Soit $\pi_0\leqslant\pi$ un cardinal qui est utile pour $\underline{S}$ et pour $\mathcal{M}$, et tel que $f:\underline{S}\longrightarrow\mathcal{M}$ soit $\pi_0$-accessible, i.e. tel que $\underline{S}$ soit stable dans $\mathcal{M}$ par $\varinjlim$ grandes devant $\pi_0$ (on peut p.ex. prendre $\pi_0=\pi$). Alors pour tout cardinal $\pi'>\pi$ tel que}
\begin{itemize}[leftmargin=2.5em]
\item[\textit{a)}] \textit{$\pi'^{\pi_0}=\pi'$ et}
\item[\textit{b)}] \textit{$\pi'\geqslant\cardred\underline{S}_{\pi_0}$,}
\end{itemize}
\textit{$\underline{S}$ est $\Phi$-$\pi'$-accessible dans $\mathcal{M}$. Donc si $\pi'>\pi$ est tel que $\pi'^{\pi}=\pi'$, $\pi'>\cardred\underline{S}_\pi$, alors $\underline{S}$ est $\Phi$-$\pi'$-accessible. (Il suffit p.ex. de prendre $\pi'=2^c$, avec $c\geqslant\sup(\pi,\cardred\underline{S}(\pi))$.) Si on admet l'HCG, il suffit donc qu'on ait $\pi'>\sup(\pi,\cardred\underline{S}(\pi))$.}

\msp{page 349}
\textbf{Proposition 4.13.28} (\footnote{À rejeter plus bas, après 4.13.31, qui est la première application importante de 4.13.22, alors que 4.13.28 se prouve sans avoir à utiliser 4.13.22 ou ses variantes.}). \textit{Soient $\mathcal{M}$, $\pi$, $\Phi$ comme dans 4.13.23. Supposons de plus que $J\in\Phi$ implique que $J$ est} filtrante. \textit{Alors l'ensemble $\tilde{\Sigma}^{\Phi}_{\pi}$ des sous-ensembles $\Phi$-$\pi$-accessibles de $\mathcal{M}$ est stable par réunions finies. Plus généralement, soit $\lambda$ un cardinal $<\pi$ tel que $\forall\ J\in\Phi$, $J$ soit grande devant $\lambda$. Alors l'ensemble $\tilde{\Sigma}^{\Phi}_{\pi}$ est stable par réunions $(S_r)_{r\in A}$, indexées par des ensembles $A$ de cardinal $\leqslant\lambda$.}

\textbf{Corollaire 4.13.29.} \textit{L'ensemble des parties accessibles de} $\Ob\mathcal{M}$ \textit{est stable par “petites” réunions -- i.e. si $(S_r)_{r\in A}$ est une petite famille de parties accessibles de} $\Ob\mathcal{M}$, \textit{alors $\bigcup_{r\in A}S_r$ est aussi accessible.}

La démonstration de 4.13.28 repose sur le lemme suivant (je n'explicite par la démonstration) : % [sic] par la démonstration -> pas la démonstration, sense of the sentence

\printedpage{133}
% head: A. Grothendieck : Dérivateurs, Ch. XVIII, Édition des Universités de Montpellier II et Paris VII
% foot: Version 17 décembre 2025
\msp{page 350}
\textbf{Lemme 4.13.30.} \textit{Soit $I$ une catégorie grande devant $\pi$, et $(J_r)_{r\in A}$ une famille des sous-catégories pleines telles que $\card A\leqslant\pi$, telle que $\bigcup_{r\in A}\Ob J_r$ soit cofinal dans $\Ord(I)$, i.e. pour tout $i\in\Ob I$ $\exists\ j\in\bigcup_{r\in A}\Ob J_r$, et une flèche $i\longrightarrow j$. Alors il existe $r\in A$ tel que $J_r$ soit cofinal dans $I$.}

\textbf{Théorème 4.13.31.} \textit{Soient $\mathcal{M}$, $\mathcal{N}$ des catégories accessibles, $f:\mathcal{M}\longrightarrow\mathcal{N}$ un foncteur accessible, $T\subset\Ob\mathcal{N}$ une partie de $\mathcal{N}$, $\underline{T}\subset\mathcal{N}$ la sous-catégorie strictement pleine correspondante de $\mathcal{N}$, $S=f^{-1}(T)$, $\underline{S}=f^{-1}(\underline{T})$ la sous-catégorie pleine de $\mathcal{M}$ associée. Si $T$ est accessible, alors $S$ est accessible\,; en d'autres termes, si $\underline{T}$ est accessible dans $\mathcal{N}$, $\underline{S}$ est accessible dans $\mathcal{M}$.}

Cela résulte de l'énoncé plus précis :

\textbf{Corollaire 4.13.32.} \textit{Soient $\pi_0$ et $\pi>\pi_0$ des cardinaux ayant les propriétés suivantes :}

\msp{page 351}
\begin{itemize}[leftmargin=2.5em]
\item[\textit{a)}] \textit{$\underline{T}$ est $\pi_0$-$\pi$-accessible dans $\mathcal{N}$ (4.13.23).}
\item[\textit{b)}] \textit{$\pi_0$ est utile pour $\mathcal{M}$ (il l'est pour $\mathcal{N}$, à cause de a)).}
\item[\textit{c)}] \textit{$f(\mathcal{M}_\pi)\subset\mathcal{N}_\pi$.}
\item[\textit{d)}] \textit{$\mathcal{M}$ satisfait à $L_{\pi_0}$, et $f$ est $\pi_0$-accessible.}
\end{itemize}
\textit{(L'existence d'un tel couple $(\pi_0,\pi)$ résulte aussitôt de 4.13.16 et de 4.12.1.) Alors $\underline{S}$ est $\pi_0$-$\pi$-accessible dans $\mathcal{N}$.}

Comme $f$ commute aux $\varinjlim$ grandes devant $\pi_0$ en vertu de d), et que $\underline{T}$ est stable par ces limites par a) (cf. 4.13.23, 3°), il en est donc de même de $\underline{S}$, i.e. le 3° de la définition 4.13.23 est satisfait. D'autre part, la condition 1° ($\pi$ utile pour $\mathcal{M}$, $\mathcal{M}$ stable par $\varinjlim$ grandes devant $\pi_0$ et de cardinal $\leqslant\pi$) est vérifiée par b) et d). Donc il reste à vérifier la condition 2°, i.e. que $\underline{S}$ est $\pi$-accessible. Posons
\[
\begin{aligned}
\underline{T}(\pi) &= \underline{T}\cap\mathcal{N}_\pi\\
\underline{S}(\pi) &= \underline{S}\cap\mathcal{M}_\pi\ \ \underbrace{\ \subset\ }_{\text{grâce à c)}}\ \ f^{-1}(\underline{T}(\pi))\ .
\end{aligned}
\]
Comme $\underline{S}$ est stable dans $\mathcal{M}$ par $\varinjlim$ grandes devant $\pi_0$ (lesquelles existent dans $\mathcal{M}$ par d), et a fortiori par $\varinjlim$ grandes devant $\pi$, on a un foncteur canonique pleinement fidèle
\[ \Ind_\pi\underline{S}(\pi)\ \longrightarrow\ \underline{S}\ . \]

\msp{page 352}
D'ailleurs, $\underline{T}$ étant strictement pleine dans $\mathcal{N}$, $\underline{S}$ est strictement pleine dans $\mathcal{M}$, et il reste à vérifier que le foncteur précédent est essentiellement surjectif. Soit donc $x\in\Ob\underline{S}$. On aura, grâce à b),
\[ x\ =\ \varinjlim_{I}x_i\ ,\qquad\quad x_i\in\mathcal{M}_\pi,\ I=\mathcal{M}_\pi/x\ . \]
D'ailleurs, le système inductif précédent est $\pi_0$-$\pi$-adapté, i.e. (p. 339) :
\begin{itemize}[leftmargin=2.5em]
\item[1)] $I$ équivalente à une petite catégorie, $I$ grande devant $\pi$ (résulte du fait que $\pi$ est utile pour $\mathcal{M}$, cf.\,\cercle{b}\,).
\end{itemize}

\end{document}
